\section{Method}\label{sec:method}


Our method addresses the instability of critic-free RLHF by replacing biased local normalization with stable global normalization. We propose two variants for different use cases.

\subsection{Global Normalization}

The primary \ours algorithm is designed for general-purpose RLHF. In its $k=1$ configuration, it is designed for maximum efficiency and prompt diversity.

It optimizes the PPO objective in \Cref{eq:ppo} but redefines the advantage $A_{q,o_t}$. We use the standard PPO reward formulation, which incorporates the k1-style KL penalty directly into the reward:
\begin{equation}\label{eq:adv}
A_{q,o_t} = r(o_{1:T},q) - \beta \cdot \sum_{i=t}^{T} \mathrm{KL}(i)
\end{equation}
where $\text{KL}(t) = \log\left(\frac{\pi_{\theta_{\text{old}}}(o_t | q, o_{<t})}{\pi_{ref}(o_t | q, o_{<t})}\right)$.

The key innovation is our normalization strategy. Instead of a non-existent (for $k=1$) or biased local norm, we use Global Advantage Normalization~\citep{andrychowicz2020matters}:
\begin{equation}\label{eq:batch_norm}
A_{q,o_t}^{\text{norm}} = \frac{A_{q,o_t} - \text{mean}\left(A|A \in \mathcal{D}_{\text{batch}}\right)}{\text{std}\left(A|A \in \mathcal{D}_{\text{batch}}\right) + \epsilon}
\end{equation}
The global normalization is the core of \ours. Pooling a large rollout batch (e.g., 1024 responses) replaces separate small-group scales with shared statistics. Under independent sampling, sample moments converge to population values (Appendix~\ref{appendix:proof_GRPO}); this is not a general unbiased-policy-gradient guarantee. Valid token positions exclude padding, and the moments remain fixed during rollout reuse.
The efficient $k=1$ case is detailed in \Cref{algo:main}.


\subsection{Local Baseline}
For more complex tasks, such as multi-step reasoning, sampling multiple responses per prompt ($k > 1$) can be beneficial \citep{yue2025vapo}. To this end, we introduce \varn. This variant combines the benefits of group-sampling with the stability of global normalization. The advantage calculation is a two-step process:
\begin{itemize}
    \item \textbf{Group Mean Subtraction for Reshaping}: We first subtract the group mean reward. It serves as a local baseline, reshaping rewards to be robust to different reward scales (e.g., $0/1$ vs. $-1/1$).
    \begin{equation}\label{eq:group_center}
     A'_{q,o_t^{(i)}} = r(q,o^{(i)}) - \frac{1}{k}\sum_{j=1}^k r(q,o^{(j)})
    \end{equation}
    \item \textbf{Global Batch Normalization for Stability}: We then normalize this initial advantage using the \textbf{global batch statistics}, not the unstable local group statistics.
    \begin{equation}
     A_{q,o_t}^{\text{norm}} = \frac{A'_{q,o_t} - \mathrm{mean}_{\text{batch}}(A')}{\text{std}_{\text{batch}}(A') + \epsilon}
    \end{equation}
\end{itemize}
This separates reward centering from scale estimation; the inclusive group baseline equals $(k-1)/k$ times RLOO (Appendix~\ref{appendix:design}).

Furthermore, this variant employs a separate $k_2$ KL regularizer. At a fixed context with on-policy actions held fixed during differentiation, its expected pointwise gradient equals the reverse-KL gradient~\citep{liu2025rethinking}. This equivalence concerns a gradient surrogate, not KL values; unweighted old-policy rollout reuse is off-policy (Appendix~\ref{appendix:kl_design}). The maximized objective is:
\begin{equation}\label{eq:baseline_loss}
\mathcal{L} = \mathcal{L}_{\text{PPO}}(A^{\text{norm}}) - \lambda\,\mathbb{E}_{o\sim\pi_{\mathrm{old}}}\left[\frac{1}{|o|}\sum_t \frac12\left(\log\frac{\pi_\theta(o_t\mid q,o_{<t})}{\pi_{ref}(o_t\mid q,o_{<t})}\right)^2\right].
\end{equation}

\subsection{Relationship with PPO}
Both variants retain PPO clipping without a learned critic. Setting $V=0$ and $\gamma=\lambda_{\mathrm{GAE}}=1$ reduces GAE to Monte Carlo returns, the structure used by the basic variant before normalization. \varn additionally uses prompt-specific reward centering and a separate KL loss; these changes distinguish it from standard actor--critic PPO.

\subsection{Implementation Tricks} \label{appendix:tricks}
\paragraph{Token-level Advantage}
The task's immediate reward is zero before termination, but its undiscounted return contributes to every preceding response token in \Cref{eq:adv}. \varn broadcasts its centered terminal reward to valid response positions. The stated objective averages token losses within each response and then across responses, as specified in Appendix~\ref{appendix:batch_mean}.

\paragraph{Batch Construction}
For REINFORCE++ ($k=1$), a global batch of size $N$ comprises $N$ distinct prompts. For \varn ($k=4$), a global batch of size $N=1024$ might contain $ 1024/4=256$ unique prompts. The global moments pool valid response positions across the rollout batch.

\paragraph{Mini-Batch Updates}
To enhance training efficiency, we implement mini-batch updates with the following characteristics:
\begin{itemize}[leftmargin=*]
    \item \textbf{Batch Processing:} Data is processed in smaller, manageable chunks rather than full-batch updates.
    \item \textbf{Multiple Updates:} Reuse detached advantages and old-policy probabilities over $E$ optimization epochs.
    \item \textbf{Stochastic Optimization:} Introduces beneficial randomness for better generalization.
\end{itemize}

\paragraph{Reward Normalization and Clipping}

We implement comprehensive reward processing to stabilize training:
\begin{itemize}[leftmargin=*]
    \item \textbf{Normalization:} Uses a shared z-score scale across the batch; standardization alone does not remove outliers.
    \item \textbf{Clipping:} Constrains reward values within predefined bounds to avoid instability.
    \item \textbf{Scaling:} Applies appropriate scaling factors for numerical stability during updates.
\end{itemize}

\subsection{Why Use Global Batch Normalization?} \label{appendix:whyglobal}

Omitting the numerical stabilizer, under independent sampling with nonzero population variance, normalization approaches the population-standardized target as $N\to\infty$ (Appendix~\ref{appendix:proof_GRPO}). This motivates global normalization in both \textsc{REINFORCE++} and \textsc{REINFORCE++}\textsubscript{\textit{w/} Baseline}: a batch of 1024 responses pools more observations than groups of 4 or 8. Token dependence and repeated responses limit the effective independent sample count. The result explains the scale estimate, rather than proving that every clipped, reused-rollout policy gradient is unbiased.

\subsection{Summary}

To address these challenges, we propose \ours, a method centered on a simple yet powerful idea: \textbf{Global Advantage Normalization}, without a critic. Instead of normalizing within a prompt's local group, \ours normalizes the advantage function across the \textit{entire global training batch}. The overall approach is:
\begin{itemize}
    \item \textbf{Theoretical Stability}: With independent sampling, increasing the batch size improves population-moment estimation. Finite-sample effects and trajectory dependence remain relevant (Appendix~\ref{appendix:proof_GRPO}).
    \item \textbf{Training Efficiency}: It retains the critic-free architecture, significantly reducing computational and memory overhead compared to PPO.
    \item \textbf{Strong Generalization}: Using a shared scale is associated with improved held-out reasoning results in our evaluated settings.
\end{itemize}

\begin{algorithm}
\caption{\ours when $k=1$}\label{algo:main}
\begin{algorithmic}[1]
\Require Initial policy model $\pi_{ref}$, reward models $R$, task prompts $\mathcal{D}$
\State policy model $\pi_{\theta} \gets \pi_{ref}$
\For {step $= 1, \ldots, M$}
    \State Sample a batch $\mathcal{D}_{batch}$ from $\mathcal{D}$
    \State Update the old policy model $\pi_{\text{old}} \gets \pi_{\theta}$
    \State \textbf{Sample output ($k=1$)} $o \sim \pi_{\text{old}}(\cdot \mid q)$ for each question $q \in \mathcal{D}_{batch}$
    \State Compute rewards $r_i$ for each sampled $o_i$ using $R$
    \State Compute advantage $A_{q,o_t}$ via \Cref{eq:adv}
    \State Normalize advantages globally across $\mathcal{D}_{batch}$ to get $A_{q,o_t}^{\text{norm}}$ via \Cref{eq:batch_norm}
    \For {iteration $= 1, \ldots, E$}
        \State Update $\pi_{\theta}$ by maximizing the objective in \Cref{eq:ppo} using $A_{q,o_t}^{\text{norm}}$
    \EndFor
\EndFor
\Ensure $\pi_{\theta}$
\end{algorithmic}
\end{algorithm}
